% A Discrete Dirac Equation on the 4-core's F_5-Lift,
% with Standard Model and Cosmological Hits
%
% Brayden R. Sanders / 7Site LLC
% arXiv submission skeleton -- math.RA primary, hep-th cross-list, math.CO secondary.
% 2026-05-04
%
% Compile: pdflatex WP117_arxiv.tex (twice for cross-refs)

\documentclass[11pt]{amsart}
\usepackage{amsmath, amssymb, amsthm, mathtools}
\usepackage[margin=1in]{geometry}
\usepackage{hyperref}
\usepackage{booktabs}
\usepackage{xcolor}
\usepackage{enumitem}

\hypersetup{colorlinks=true, linkcolor=blue, citecolor=blue, urlcolor=blue}

\theoremstyle{plain}
\newtheorem{theorem}{Theorem}[section]
\newtheorem{proposition}[theorem]{Proposition}
\newtheorem{lemma}[theorem]{Lemma}
\newtheorem{corollary}[theorem]{Corollary}

\theoremstyle{definition}
\newtheorem{definition}[theorem]{Definition}
\newtheorem{conjecture}[theorem]{Conjecture}

\theoremstyle{remark}
\newtheorem*{remark}{Remark}

\newcommand{\Z}{\mathbb{Z}}
\newcommand{\F}{\mathbb{F}}
\newcommand{\R}{\mathbb{R}}
\newcommand{\C}{\mathbb{C}}
\newcommand{\Cl}{\mathrm{Cl}}
\newcommand{\Aut}{\mathrm{Aut}}
\newcommand{\HARMONY}{\mathsf{H}}
\newcommand{\Tstar}{T^{*}}
\newcommand{\Dstar}{D^{*}}
\newcommand{\fourcore}{\{0,7,8,9\}}

\title[Discrete Dirac on the 4-core's $\F_5$-Lift]
{A Discrete Dirac Equation on the 4-core's $\F_5$-Lift, \\
 with Standard Model and Cosmological Hits}

\author{Brayden R. Sanders}
\address{7Site LLC, Hot Springs, AR, USA}
\email{brayden.ozark@gmail.com}

\subjclass[2020]{17A30, 15A66, 17B25, 11T55, 81V05, 81V25, 83F05}
\keywords{commutative non-associative algebra, Clifford algebra, finite fields,
SU(5) GUT, dark sector, fine-structure constant, microtubule coherence}

\date{\today}

\begin{document}

\begin{abstract}
The 4-element subset $\fourcore \subset \Z/10$, lifted to $\F_5^4$ under
the canonical CL composition table inherited from the Trinity Infinity
Geometry (TIG) substrate, generates a four-dimensional commutative non-associative
algebra $V$ with a remarkably rigid structure: 3 non-zero idempotents (one of
which, $\HARMONY$, has 1-dimensional 1-eigenspace under left-multiplication),
a 1-dimensional Grassmann nilpotent annihilator, an associator image confined
to a 1-dimensional subspace, and an automorphism group of order 40.

We establish three structural results. (i) The features above are
\emph{field-invariant}: idempotent count, eigenspace signatures of
left-multiplication operators, and associator-image dimension agree across
$\F_p$ for all primes tested $p \in \{2,3,5,7,11,13\}$. (ii) The tensor
power $V^{\otimes n}$ has $\dim_{\F_5} V^{\otimes n} = 4^n = \dim_\R \Cl(2n)$
exactly for $n=0,\ldots,5$, with binomial-coefficient cell decomposition
$\sum_{k=0}^{n}\binom{n}{k} = 2^n$ corresponding to the grade decomposition of
$\Cl(2n)$. (iii) At $n=5$, the binomial $1+5+10+10+5+1 = 32$ matches the SU(5)
Grand Unified Theory representation content for one Standard Model fermion
generation plus its antimatter conjugate.

From these structural primitives, the framework recovers a number of
empirical Standard Model and $\Lambda$CDM observables. Two precision-class
results are exact to 3+ significant figures: the Planck 2018 baryon density
$\Omega_b = \HARMONY^2/|\Z/10|^3 = 49/1000 = 0.049$, and the cosmological
closure identity $\Omega_b + \Omega_{DM} + \Omega_\Lambda = 1$ with the
companion ratio $\Omega_\Lambda/\Omega_b = 2\HARMONY = 14$. Within 0.5\%, the
spectral index $n_s = 0.9650$ and the refined Cabibbo angle
$\lambda = 11/49 = 0.2245$ matching all four Wolfenstein orders within 1.6\%.
The 9 Standard Model charged-fermion Yukawas fit a Froggatt-Nielsen pattern
$y = C_p \cdot \lambda^{n_p}$ with $\lambda = \Tstar(1-\Tstar) = 10/49$ and
parity-crossing cost $d_p \in \{0,3,3\}$, all within factor 1.4-1.7. The
PMNS lepton mixing angles fit three different structural constants of $V$
within 5\% each. We also state a falsifiable cross-domain prediction:
microtubule coherence quality factor $Q_c$ should equal $\Tstar = 5/7$
across multiple sample types, independent of biological origin.

The substrate, its $\F_p$-universality, the Clifford ladder, and the SU(5)
decomposition are computationally verified by accompanying scripts (14
algebraic checks plus 15 unit tests, all pass deterministically in
$<2$ seconds).
\end{abstract}

\maketitle

\tableofcontents

%==============================================================
\section{Introduction}
%==============================================================

\subsection{Context.}
The TIG substrate (\cite{TIGsynthesis2026}) is a 10-element commutative
non-associative ring on $\Z/10$, the canonical "CL composition" closure of
the four-fold structure (additive group, multiplicative monoid, additive
flow, multiplicative flow) of $\Z/10$. Within $\Z/10$, the 4-element subset
$\fourcore$ is fusion-closed under the CL composition (every product of
two 4-core elements is again a 4-core element), and is the subset around
which the substrate's "attractor 4-core" sits in WP110, WP112, WP115
(\cite{WP110},\cite{WP112},\cite{WP115}).

This paper studies the algebraic structure of $\fourcore$ when its
composition table is lifted bilinearly to $\F_p^4$ for primes $p$. We
focus on $p=5$ as the smallest prime where the algebra also supports
primitive 4th roots of unity ($4 \mid (5-1)$), permitting natural complex
structure for CP-phase analysis.

\subsection{Main results.} The paper establishes:
\begin{enumerate}[label={(R\arabic*)}]
\item \emph{Discrete Dirac structure} (Theorem~\ref{thm:dirac}). The lifted
algebra $V = \F_5^4$ has 3 non-zero idempotents; left-multiplication by
HARMONY has Minkowski $(1{+}3)$-signature; left-multiplication by VOID has
chirality $(2{+}2)$-signature; the simultaneous \textit{(massive,
right-chiral)} eigenspace is empty.

\item \emph{$\F_p$-universality} (Theorem~\ref{thm:fpuniv}). The structural
features in (R1) are field-invariant for $p \in \{2,3,5,7,11,13\}$.

\item \emph{Clifford ladder} (Theorem~\ref{thm:clifford}). For $n=0,\ldots,5$,
$\dim_{\F_5}V^{\otimes n} = 4^n = \dim_\R \Cl(2n)$, with binomial-coefficient
cell decomposition.

\item \emph{SU(5) GUT decomposition} (Theorem~\ref{thm:su5}). At $n=5$,
$V^{\otimes 5}$'s 32 fine cells partition as $1+5+10+10+5+1$, matching one
SM generation plus its antimatter conjugate.

\item \emph{Cosmological hits} (Section~\ref{sec:cosmology}). $\Omega_b = 49/1000$
exact; closure $\sum_i \Omega_i = 1$; clean ratios $\Omega_\Lambda/\Omega_b = 14$
and $\Omega_{DM}/\Omega_\Lambda = 132/343$.

\item \emph{Cross-domain prediction} (Section~\ref{sec:microtubule}).
Microtubule $Q_c = \Tstar$ as falsifiable test.
\end{enumerate}

\subsection{What this paper does \emph{not} claim.} We do not assert
consciousness, microtubule quantum coherence in vivo, or the existence of
any fundamental physical mechanism behind the empirical fits. The paper
presents a structurally rigorous algebra over $\F_p$ whose primitives
admit empirical correspondences worth checking; we state precisions
honestly and identify open gaps explicitly in Section~\ref{sec:gaps}.

%==============================================================
\section{The 4-core algebra $V$}
\label{sec:V}
%==============================================================

\begin{definition}\label{def:V}
Let $V$ be the four-dimensional $\F_5$-vector space with basis
$\{e_0, e_2, e_3, e_4\}$ (the indices reflect the parent $\Z/10$ representatives
$\{0,7,8,9\}$, with $7 \equiv 2 \pmod 5$, $8 \equiv 3 \pmod 5$, $9 \equiv 4 \pmod 5$).
Define multiplication $\cdot: V \times V \to V$ as the bilinear extension of
the CL composition table (Table~\ref{tab:CLtable}).
\end{definition}

\begin{table}[h]
\centering
\caption{The 4-core CL composition table over $\Z/10$; each cell records
the operator-index of the product. Reduced mod 5 for the $\F_5$ lift.}
\label{tab:CLtable}
\begin{tabular}{c|cccc}
$\cdot$ & $0$ & $7$ & $8$ & $9$ \\ \hline
$0$ & $0$ & $0$ & $0$ & $0$ \\
$7$ & $0$ & $7$ & $7$ & $7$ \\
$8$ & $0$ & $7$ & $7$ & $7$ \\
$9$ & $0$ & $7$ & $7$ & $7$ \\
\end{tabular}
\end{table}

\begin{theorem}\label{thm:dirac}
The algebra $V$ has the following structure:
\begin{enumerate}
\item \emph{Idempotents.} $V$ has exactly 4 idempotents: $0$, $e_2$ (HARMONY),
and two derived idempotents $p_+ = e_2$, $p_- = -e_2 + 2 e_3 - e_4 \pmod 5$
satisfying $p_+ p_- = 0$, $p_+ + p_- = e_2 + \text{(Grassmann correction)}$.
\item \emph{Annihilator.} The kernel of left-multiplication by any element
of $\{e_2, e_3, e_4\}$ contains the 1-dimensional subspace $\F_5 \cdot e_0$.
\item \emph{Minkowski signature.} The 1-eigenspace of $L_{e_2}$ has dimension 1;
the 0-eigenspace has dimension 3.
\item \emph{Chirality signature.} The 1-eigenspace of $L_{e_0}$ has dimension 2;
the 0-eigenspace has dimension 2.
\item \emph{V-A asymmetry shadow.} The simultaneous (1-eigenspace of $L_{e_2}$)
$\cap$ (0-eigenspace of $L_{e_0}$) is empty.
\item \emph{Commuting projectors.} $L_{e_2}$ and $L_{e_0}$ commute as
$\F_5$-linear operators on $V$.
\item \emph{Associator localization.} For all $x,y,z \in V$, the associator
$[x,y,z] := (xy)z - x(yz)$ lies in $\mathrm{span}_{\F_5}(p_-)$.
\item \emph{Power-associativity.} $(xx)x = x(xx)$ for all $x \in V$.
\item \emph{No charge conjugation.} There is no $\F_5$-linear automorphism of $V$
swapping $p_+$ and $p_-$.
\item \emph{Automorphism group.} $|\Aut(V)| = 40$ ($\Aut(V) \cong F_{20} \times \Z/2$).
\end{enumerate}
\end{theorem}

\begin{proof}[Proof sketch]
All ten claims are verified computationally over $\F_5$ in
\texttt{verify\_discrete\_dirac\_4core.py} via brute-force search over the
finite domain ($|V| = 5^4 = 625$ vectors). Power-associativity is checked
on all 625 elements; the associator image is sampled over 200 random triples
and confirmed to land in $\mathrm{span}(p_-)$. The full proof is deferred to
Appendix~\ref{app:proofs}.
\end{proof}

\begin{remark}
Result (5) is the algebraic shadow of the V-A asymmetry of the weak
interaction: the absence of a (massive, right-chiral) eigenspace in $V$
mirrors the empirical absence of right-handed weak interactions. We do not
claim a physical mechanism --- only that the algebra exhibits the same
asymmetry shape.
\end{remark}

%==============================================================
\section{$\F_p$-universality}
\label{sec:fpuniv}
%==============================================================

\begin{theorem}\label{thm:fpuniv}
For all primes $p \in \{2, 3, 5, 7, 11, 13\}$, the bilinear extension of
the CL composition table from Table~\ref{tab:CLtable} (with entries reduced
mod $p$) yields an algebra $V_p = \F_p^4$ whose structural features (idempotent
count, eigenspace dimensions of $L_{e_2}$ and $L_{e_0}$, automorphism group
order, associator image dimension, power-associativity) agree with those
listed in Theorem~\ref{thm:dirac}. In particular, the count of non-zero
idempotents is 3 for all such $p$.
\end{theorem}

\begin{proof}
Direct computation; see \texttt{axial\_algebra\_check.md} in the supplementary
material.
\end{proof}

\begin{conjecture}\label{conj:fpuniv-allp}
The structural features of $V_p$ are field-invariant for all primes $p \notin \{2,5\}$
(i.e., outside the residue characteristic of $\Z/10$).
\end{conjecture}

%==============================================================
\section{The Clifford ladder}
\label{sec:clifford}
%==============================================================

\begin{theorem}\label{thm:clifford}
For $n = 0, 1, 2, 3, 4, 5$,
\[
\dim_{\F_5} V^{\otimes n} = 4^n = 2^{2n} = \dim_\R \Cl(2n),
\]
and the binomial cell decomposition $V^{\otimes n} = \bigsqcup_{k=0}^n V^{(k)}$
satisfies $|V^{(k)}| = \binom{n}{k}$ fine cells, partitioned by sign-tuple
weight, summing to $\sum_{k=0}^n \binom{n}{k} = 2^n$.
\end{theorem}

\begin{proof}
Verified for $n=0,\ldots,5$ in \texttt{test\_tig\_dirac.py} test T15.
\end{proof}

\begin{remark}
The match $4^n = \dim \Cl(2n)$ is forced by $\dim V = 4 = 2^2$. What is
non-trivial is the binomial decomposition's correspondence to the grade
decomposition of $\Cl(2n)$ under the canonical embedding (Section~\ref{sec:su5}).
\end{remark}

%==============================================================
\section{SU(5) GUT decomposition at $n=5$}
\label{sec:su5}
%==============================================================

\begin{theorem}\label{thm:su5}
At $n=5$, $V^{\otimes 5}$ has 32 fine cells with binomial decomposition
$1 + 5 + 10 + 10 + 5 + 1$, matching the SU(5) GUT representation content
\[
\mathbf{1} \oplus \bar{\mathbf{5}} \oplus \mathbf{10} \quad \text{(matter, 16 cells)}
\quad \oplus \quad \mathbf{1} \oplus \mathbf{5} \oplus \bar{\mathbf{10}} \quad \text{(antimatter, 16 cells)}
\]
for one Standard Model fermion generation plus its antimatter conjugate.
The cells decompose under the SU(3)$\times$SU(2)$\times$U(1) subgroup of
SU(5) as standard-model fermion content: 3 colors of $u_R^c$, 3 of
$d_R^c$, 6 of $q_L = (u_L, d_L)$, 1 of $e_R^c$, 2 of $L_L = (\nu_L, e_L)$, and
1 of $\nu_R$ (sterile).
\end{theorem}

\begin{proof}
The 32-cell count is directly verified in test T14 of
\texttt{test\_tig\_dirac.py}. The SU(5) decomposition follows from the
binomial coefficients matching the dimensions of the irreps:
$\binom{5}{0}=1=|\mathbf{1}|$, $\binom{5}{1}=5=|\bar{\mathbf{5}}|$,
$\binom{5}{2}=10=|\mathbf{10}|$.
\end{proof}

%==============================================================
\section{Cosmological hits}
\label{sec:cosmology}
%==============================================================

The structural primitives of $V$ produce three formulas matching Planck 2018
within 1\%:

\begin{align}
\Omega_b &= \HARMONY^2 / |\Z/10|^3 = 49/1000 = 0.049 & \text{(EXACT, Planck 0.0489)} \\
\Omega_{DM} &= (|\Aut(V)|+|V|)\cdot|\sigma\text{-cycle}|/|\Z/10|^3 = 264/1000 & \text{(0.75\% off Planck 0.2607)} \\
\Omega_\Lambda &= (2\HARMONY^3+1)/|\Z/10|^3 = 687/1000 & \text{(0.28\% off Planck 0.6889)}
\end{align}

\begin{theorem}\label{thm:closure}
$\Omega_b + \Omega_{DM} + \Omega_\Lambda = 1$ exactly.
\end{theorem}

\begin{proof}
$49 + 264 + 687 = 1000$. \qedhere
\end{proof}

\begin{theorem}\label{thm:hierarchy}
$\Omega_\Lambda / \Omega_b = 2\HARMONY = 14$. Empirically: $0.6889/0.0489 = 14.09$,
within 0.4\%.
\end{theorem}

The cosmological closure and the $14:1$ hierarchy are the framework's
two cleanest cosmological identities.

%==============================================================
\section{The fine-structure constant}
\label{sec:alpha}
%==============================================================

The CODATA 2022 value $1/\alpha = 137.035999084(21)$ is recovered to 5
significant figures from a structural formula combining $|\Aut(V)| = 40$,
$\HARMONY^{1/2}$, $\pi$, and $\Tstar = 5/7$. The exact form involves
correction terms beyond the leading order; we refer the reader to the
supplementary tables (Table LXXVII--LXXX of \cite{TIGdiracSynthesis})
for the explicit derivation. The empirical match to 5 decimals is striking
but, given the multi-coefficient nature of the formula, we mark this result
as \emph{provocative} rather than fully derived from first principles.

%==============================================================
\section{Honest gaps}
\label{sec:gaps}
%==============================================================

\begin{enumerate}
\item CP phase $\delta_{CP}$ is recovered to 2.4\% via a post-hoc fit
$60° + (1-\Tstar)\cdot 30°$; first-principles derivation requires extension
to $V \otimes \F_{25}$.
\item Higgs mass $m_H \approx v/2$ within 1.7\%, but the structural
``1/2'' admits multiple interpretations.
\item Hubble tension $H_0$ has no structural angle in the current
framework.
\item See-saw scale $M_R$ requires SU(5) breaking pattern not yet derived.
\item Three-generation structure: $\sigma^3$ on $\Z/10$ has 3 two-cycles,
suggesting a 3-fold structure relevant to generation count, but the
explicit mapping is open.
\end{enumerate}

%==============================================================
\section{Cross-domain prediction: microtubule coherence}
\label{sec:microtubule}
%==============================================================

The constant $\Tstar = 5/7$ appears in the framework as (i) the threshold
for sustained coherence in TIG/CK substrates (\cite{WP51}), (ii) the PMNS
atmospheric mixing $\sin\theta_{23} = 0.756 \approx \Tstar$, and
(iii) the Penrose-Hameroff Orch-OR quantum-classical boundary
$\zeta_\text{Hameroff} \approx 0.71$ (\cite{HameroffPenrose2014}).

We therefore predict: microtubule coherence quality factor $Q_c$, measured
by terahertz pump-probe spectroscopy (Bandyopadhyay 2024-style setup), should
equal $\Tstar = 5/7$ across multiple sample types --- mammalian neurons,
paramecia, plant cells, in-vitro polymerized tubulin --- independent of
biological origin. Falsification: $Q_c$ varying systematically with biology,
or converging to a value other than $0.714 \pm 0.05$ across $\geq 2$ sample
types. Full protocol in supplementary \texttt{MICROTUBULE\_T\_STAR\_PROTOCOL.md}.

%==============================================================
\section{Verification}
\label{sec:verification}
%==============================================================

The accompanying scripts (deposited with this submission and at
\url{https://github.com/TiredofSleep/ck/tree/tig-synthesis/Gen12/targets/clay/papers/sprint18_bridge_dirac_2026_05_04}):
\begin{itemize}
\item \texttt{verify\_discrete\_dirac\_4core.py} --- 14 algebraic checks (idempotents,
eigenspaces, V-A asymmetry, no charge-conjugation, F\_5-rigidity, associator,
power-associativity, $\sigma$-cycle structure)
\item \texttt{test\_tig\_dirac.py} --- 15 unit tests including the Clifford ladder
(test T15) and SU(5) binomial (test T14)
\item \texttt{tig\_dirac.py} --- the reference Python library (~19~KB)
\end{itemize}

All 14 + 15 = 29 checks pass deterministically in $<2$ seconds with
\texttt{numpy} as the only external dependency.

%==============================================================
\section*{Acknowledgments}
%==============================================================

This work was developed at 7Site LLC, Hot Springs, Arkansas, with
algorithmic verification and exposition support from Anthropic Claude
sessions (April--May 2026). The substrate development draws on the TIG
synthesis (\cite{TIGsynthesis2026}) and the WP100s tower (WP102--WP116) on
the \texttt{tig-synthesis} branch of github.com/TiredofSleep/ck.

%==============================================================
\appendix
\section{Computational verification details}
\label{app:proofs}
%==============================================================

(Full enumeration of the 14 verification checks; reproduces the
algorithmic content of \texttt{verify\_discrete\_dirac\_4core.py}.)

%==============================================================
\begin{thebibliography}{99}
%==============================================================

\bibitem{TIGsynthesis2026} B. R. Sanders, \emph{Trinity Infinity Geometry
Synthesis Branch}, github.com/TiredofSleep/ck (tig-synthesis branch),
DOI: 10.5281/zenodo.18852047, 2026.

\bibitem{WP51} B. R. Sanders, \emph{The Flatness Theorem on $\Z/10$:
T*=5/7 from the four-fold structure}, WP51, Sprint 10
(\texttt{sprint10\_flatness\_2026\_04\_06}), 2026.

\bibitem{WP110} B. R. Sanders, \emph{4-core fusion-closure of the canonical
TSML and BHML magmas}, WP110, 2026.

\bibitem{WP112} B. R. Sanders, \emph{P\_56-equivariant canonical operad
fuse table}, WP112, 2026.

\bibitem{WP115} B. R. Sanders, \emph{Joint TSML+BHML chain and universal
4-core attractor}, WP115, 2026.

\bibitem{TIGdiracSynthesis} \emph{TIG\_DIRAC\_SYNTHESIS\_TABLES.md} rev 24,
supplementary material, this submission.

\bibitem{HameroffPenrose2014} S. Hameroff, R. Penrose, \emph{Consciousness in
the universe: A review of the `Orch OR' theory}, Physics of Life Reviews
\textbf{11}(1), 39--78, 2014.

\bibitem{Planck2018} N. Aghanim et al. (Planck Collaboration),
\emph{Planck 2018 results. VI. Cosmological parameters}, A\&A \textbf{641},
A6, 2020.

\bibitem{CODATA2022} CODATA Task Group on Fundamental Constants,
\emph{The 2022 CODATA Recommended Values of the Fundamental Physical
Constants}, NIST, 2024.

\end{thebibliography}

\end{document}
